% ======================================================================
\chapter{A conta no GitHub}
\label{cap:github}
% ======================================================================

\section{O que é o GitHub}

\begin{conceito}[Git e GitHub não são a mesma coisa]
O \textbf{Git} é um programa de controle de versão que roda no seu computador.
O \textbf{GitHub} é uma plataforma na internet que hospeda repositórios Git e
acrescenta o que uma equipe precisa para colaborar: revisão de código (Pull
Requests), controle de acesso, Issues, publicação de sites (GitHub Pages) e
automação (GitHub Actions).
\end{conceito}

O Git funciona sozinho, sem internet. O GitHub entra quando é preciso
compartilhar o trabalho, revisar o de outras pessoas ou publicar algo.

\section{Criando a conta}

\begin{enumerate}
  \item Acesse \url{https://github.com} e clique em \menu{Sign up}.
  \item Informe e-mail, senha e nome de usuário
        (Figura~\ref{fig:signup_page}). O nome de usuário aparece em todos os
        seus endereços, como \texttt{github.com/seu-usuario}: escolha um que
        você use profissionalmente.
  \begin{figure}[H]
    \centering
    \includegraphics[width=0.8\textwidth]{./assets/images/01_signup_page.png}
    \caption{Página de criação de conta no GitHub}
    \label{fig:signup_page}
  \end{figure}
  \item Confirme o e-mail pelo código enviado e conclua as telas seguintes. Ao
        final, você verá a página inicial (Figura~\ref{fig:signed_page}).
  \begin{figure}[H]
    \centering
    \includegraphics[width=0.8\textwidth]{./assets/images/02_signed_page.png}
    \caption{Primeira visão do GitHub depois de criar a conta}
    \label{fig:signed_page}
  \end{figure}
\end{enumerate}

\section{Autenticação em dois fatores (2FA)}

O GitHub exige a autenticação em dois fatores de quem contribui com código na
plataforma. Quando a exigência chega à sua conta, você recebe um aviso e tem
45 dias para ativar; depois disso, há mais 7 dias de carência antes de o
acesso ser bloqueado \cite{github_2fa}. Como todo participante desta
capacitação vai publicar código, vale ativar logo no primeiro dia.

\begin{enumerate}
  \item Clique na sua foto (canto superior direito) \(\rightarrow\)
        \menu{Settings} \(\rightarrow\) \menu{Password and authentication}.
  \item Em \menu{Two-factor authentication}, clique em
        \menu{Enable two-factor authentication}.
  \item Leia o QR code com um aplicativo autenticador no celular e digite o
        código de seis dígitos que ele mostrar.
  \item \textbf{Baixe e guarde os códigos de recuperação}. Eles são a única
        forma de entrar na conta se você perder o celular.
\end{enumerate}

\begin{dica}[Método principal e reserva]
O GitHub recomenda um \textbf{aplicativo autenticador} (TOTP) como método
principal e uma \emph{passkey} ou chave de segurança como reserva. Passkeys
ainda não são aceitas como método principal.
\end{dica}

\section{Benefícios para estudantes}

\subsection{O que o GitHub Education oferece}

Estudantes verificados pelo GitHub Education recebem, sem custo e enquanto
forem estudantes:

\begin{itemize}
  \item o plano \textbf{GitHub Pro} na conta pessoal;
  \item o \textbf{Copilot Student}, um plano do GitHub Copilot próprio para
        estudantes, ativado em um passo separado \cite{github_copilot_student};
  \item o \textbf{Student Developer Pack}, com ofertas de ferramentas de
        empresas parceiras.
\end{itemize}

\begin{table}[H]
  \centering
  \caption{Diferenças entre o GitHub Free e o GitHub Pro (conta pessoal)}
  \label{tab:free_pro}
  \begin{tabularx}{\textwidth}{@{}X c c@{}}
    \toprule
    \textbf{Recurso} & \textbf{GitHub Free} & \textbf{GitHub Pro} \\
    \midrule
    Repositórios públicos e privados & ilimitados & ilimitados \\
    Minutos de GitHub Actions por mês\textsuperscript{a} & 2.000 & 3.000 \\
    Armazenamento do GitHub Packages & 500 MB & 2 GB \\
    Codespaces: horas de núcleo por mês & 120 & 180 \\
    Codespaces: armazenamento por mês & 15 GB & 20 GB \\
    Git LFS: armazenamento e banda por mês & 10 GiB & 10 GiB \\
    GitHub Pages em repositório privado & \nok & \ok \\
    Revisores obrigatórios em repositório privado & \nok & \ok \\
    Suporte & comunidade & e-mail \\
    \bottomrule
  \end{tabularx}
  \par\smallskip
  {\footnotesize \textsuperscript{a} Os minutos só são contados em
  repositórios privados: em repositórios públicos, o GitHub Actions é gratuito.}
\end{table}

\begin{dica}[Copilot também existe no plano gratuito]
Qualquer conta pode usar o \textbf{Copilot Free}, com limites mensais (por
exemplo, 2.000 sugestões de código por mês). O Copilot Student libera mais
recursos para quem é estudante verificado \cite{github_copilot_plans}.
\end{dica}

\subsection{Quem pode pedir}

Segundo a documentação oficial \cite{github_education_apply}, é preciso:

\begin{itemize}
  \item estar matriculado em um curso que confere diploma (escola, faculdade,
        universidade);
  \item ter pelo menos 13 anos;
  \item ter uma conta pessoal no GitHub;
  \item enviar um comprovante de matrícula atual: foto da carteira de
        estudante com a data de validade, grade de horários, histórico ou
        declaração de matrícula.
\end{itemize}

\subsection{Passo a passo}

\begin{enumerate}
  \item \textbf{Cadastre o e-mail acadêmico.} Em \menu{Settings}
        \(\rightarrow\) \menu{Emails}, adicione o e-mail da instituição
        (Figura~\ref{fig:adding_academic_email}) e confirme pelo link que
        chegar na caixa de entrada.
  \begin{figure}[H]
    \centering
    \includegraphics[width=0.9\textwidth]{./assets/images/03_adding_academic_email.png}
    \caption{Adicionando o e-mail acadêmico}
    \label{fig:adding_academic_email}
  \end{figure}
  \item \textbf{Abra os benefícios de educação.} Acesse
        \url{https://github.com/settings/education/benefits} e, em
        \menu{GitHub Education}, clique em \menu{Start an application}
        (Figura~\ref{fig:github_student_developer_pack}). A página pública do
        pack (Figura~\ref{fig:github_student}) leva ao mesmo formulário.
  \begin{figure}[H]
    \centering
    \includegraphics[width=0.9\textwidth]{./assets/images/04_gsdp_page.png}
    \caption{Página do Student Developer Pack (\texttt{education.github.com/pack})}
    \label{fig:github_student}
  \end{figure}
  \begin{figure}[H]
    \centering
    \includegraphics[width=0.9\textwidth]{./assets/images/05_application_page.png}
    \caption{\menu{Education benefits} nas configurações, com o botão \menu{Start an application}}
    \label{fig:github_student_developer_pack}
  \end{figure}
  \begin{figure}[H]
    \centering
    \includegraphics[width=0.9\textwidth]{./assets/images/06_start_application.png}
    \caption{Formulário da candidatura}
    \label{fig:application_start}
  \end{figure}
  \item \textbf{Preencha e envie.} Escolha \menu{Student}, informe a
        instituição e o e-mail acadêmico (Figura~\ref{fig:application_start}),
        anexe o comprovante
        e clique em \menu{Submit application}. A análise pode levar alguns
        dias.
  \item \textbf{Ative o Copilot Student depois da aprovação.} A aprovação e a
        ativação do Copilot são passos separados
        \cite{github_copilot_student}. Volte a
        \url{https://github.com/settings/education/benefits} e siga as
        instruções de ativação. Se só aparecerem planos pagos, espere alguns
        dias: o benefício pode demorar para ser aplicado.
\end{enumerate}

\begin{dica}[Experiência do autor]
A foto da carteira de estudante, com nome, foto e data de validade visíveis,
costuma ser analisada mais rápido que outros documentos.
\end{dica}

\begin{aviso}[As telas mudam]
As capturas deste capítulo são de 2025. O GitHub reorganiza as telas com
frequência; se algo estiver diferente, procure pelos nomes dos botões citados
no texto.
\end{aviso}
